\documentclass{article}

\usepackage[spanish]{babel}

% Required for inserting images
\usepackage{graphicx}

\usepackage[utf8]{inputenc}
\usepackage{glossaries}
\usepackage{csquotes}
\usepackage{parskip}

%Imports biblatex package
\usepackage{biblatex}

\usepackage{hyperref}
\hypersetup{
    colorlinks,
    citecolor=black,
    filecolor=black,
    linkcolor=black,
    urlcolor=black
}

% import the bibliography file
\addbibresource{investigacion.bib}

\setlength{\parindent}{0pt}

\title{Estrategias de aleatoreidad para secuenciadores electrónicos}
\author{Aarón Montoya-Moraga}

\date{diciembre 2025}

\begin{document}

\maketitle

\renewcommand*\contentsname{Tabla de contenidos}

\tableofcontents

\clearpage

\section{Resumen}


\section{Palabras clave}

\clearpage

\section{Introducción}

Este documento explica distintas estrategias usadas para aleatoreizar secuencias en secuenciadores, con un fuerte énfasis en instrumentos de fuente abierta.



\clearpage

\section{Sobre secuenciadores}



\section{Conclusiones y pasos a seguir}



\cite{korgBerlin}

\clearpage

\printbibliography[title={Bibliografía}, heading=bibintoc]

\end{document}